\section{Introduction}
\textsc{Phi} is a constructed language. A constructed language is one whose
phonology, grammar, and vocabulary, instead of having developed naturally, are
consciously devised.\footnote{\url{https://en.wikipedia.org/wiki/Constructed%
\_language}} The sound of the name \emph{Phi} embodies the simplicity and
softness of the language itself. \emph{Phi~($\varPhi$)} is also beauty in
mathematics. It is

\begin{quote}
  \emph{the golden ratio, an irrational number with a value of approximately
  1.618033988 which expresses the relationship that the sum of two quantities
  is to the larger quantity as the larger is to the smaller.} -- Wiktionary%
  \footnote{\url{https://en.wiktionary.org/wiki/\%CF\%86}}
\end{quote}

It is thought that the value of \emph{phi} is one of the fundamental
characteristics of nature; in the language Phi, nature is fundamental to
its being.

There is also an emphasis on simplicity in the language. Phi has many aspects of
a pidgin language without having been derived from existing languages. An a
priori pidgin is definitely a contradiction in terms, but the resemblance to a
pidgin is intentional. Wikipedia describes some of the attributes of a pidgin
language:\footnote{\url{https://en.wikipedia.org/wiki/Pidgin}}
\subsection{Similarities between Phi and a pidgin}
\begin{itemize}[label=\raisebox{0.33ex}{\tiny$\bullet$}]\itemsep-5pt
  \item Phi resembles an isolating language;
  \item Phi is also an analytic language;
  \item There are no embedded clauses in Phi;
  \item Phi eschews syllable codas;
  \item Consonant clusters are restricted to voiceless digraphs, e.g. \dscopy /sh/\normalfont;
  \item Aspiration only occurs with initial \dscopy /p/~\normalfont and \dscopy /t/\normalfont;
  \item Phi uses the basic five-vowel system: \dscopy /i/ /u/ /e/ /o/ /a/\normalfont;
  \item There is no sound influence between morphemes;
  \item There are no tones in Phi;
  \item No grammatical tense occurs in Phi;
  \item No conjugation or declension;
  \item There is a definite lack of grammatical gender or number;
\end{itemize}
\subsection{Difference between Phi and a pidgin}
\begin{itemize}[label=\raisebox{0.33ex}{\tiny$\bullet$}]\itemsep-5pt
  \item There are no historical sound changes in an a priori language;
  \item No monopthongization occurs in Phi; each vowel is pronounced separately;
  \item Phi has clear parts of speech and word categorization;
\end{itemize}